%%%%%%%%%%%%%%%%%%%%%%%%%%%%%%%%%%%%%%%%%%%%%%%%%%%%%%%%%%%%%
% ==================== 第五章：模型的建立与求解 ====================
% BZD数模社制作，本文档为论文模板；如需详细说明，可见论文模板完整版。
% ✓ 使用说明：
%   □ 四个问题均按“思路—建立—求解—结果—检验”组织。
%   □ 数据处理不是固定小节；只有模型确实需要时，才在建模思路或求解步骤中简述。
%   □ 模型建立末尾必须汇总最终模型，不能只分散给出若干公式。
%   □ 模型求解使用 Step 1、Step 2……说明真实求解流程，步骤数量按实际算法调整。
%   □ 结果与分析只引用真实计算结果，禁止预置模型、指标或数值。
%   □ 模型检验与灵敏度分析只选择与本问匹配的 1～2 项，不机械堆砌。
%   □ 预测类可检验误差或残差；评价类可检验排序或权重稳定性；
%     优化类可检验约束可行性、基准对比或重复求解稳定性；
%     仿真与机理类可检验边界条件、守恒关系或参数扰动。
%   □ 数组越界、维度一致、程序成功运行等属于代码自检，不写成模型检验。

\section{模型的建立与求解}

%%%%%%%%%%%%%%%%%%%%%%%%%%%%%%%%%%%%%%%%%%%%%%%%%%%%%%%%%%%%%
% ==================== 5.1：问题一的模型建立与求解 ====================
% ✓ 自查：模型选择有依据；模型已汇总；步骤可复现；结果真实；检验与模型匹配。

\subsection{问题一的模型建立与求解}

\subsubsection{建模思路与模型选择}

问题一要求【概括问题一的任务与目标】，其核心是回答【需要解决的关键数学问题】。
根据【题目条件、研究对象或数据特征】，该问题属于【数学问题类型】。
综合比较【候选方法】在【适用条件、解释能力或求解效率】方面的特点，本文选择
【主模型名称】刻画【变量关系或现实机理】，并以【必要的对比方法】作为参照。
若本问确有数据处理需求，则采用【处理方法】处理【缺失、异常、量纲或口径问题】，理由是【选择依据】；
所得【参数、特征或中间结果】将作为问题二的输入。

\subsubsection{模型的建立}

定义【决策变量、状态变量或解释变量】为【符号及含义】，定义【关键参数】为【符号及含义】。
依据【理论、机理或题目规则】建立【核心方程、目标函数或判别规则】，并结合
【资源限制、边界条件或逻辑关系】给出相应约束：

\begin{equation}
  \text{【核心方程、目标函数或判别规则】},
  \label{eq:q1_core}
\end{equation}

\begin{equation}
  \text{【约束条件、边界条件或适用条件】}.
  \label{eq:q1_constraint}
\end{equation}

\noindent\textbf{模型汇总：}
将上述变量关系、约束条件和参数范围统一整理，得到问题一的完整模型

\begin{equation}
  \mathcal{M}_1:\quad
  \left\{
  \begin{aligned}
    &\text{【核心模型表达式】},\\
    &\text{【约束或适用条件】},\\
    &\text{【变量范围、初始条件或边界条件】}.
  \end{aligned}
  \right.
  \label{eq:q1_summary}
\end{equation}

\subsubsection{模型的求解}

\noindent\textbf{Step 1：}整理【原始输入、题设参数或初始条件】，得到【可直接用于求解的输入形式】。

\noindent\textbf{Step 2：}按照【算法或计算方法】计算【关键变量、中间量或候选解】，并依据【规则】完成【核心求解过程】。

\noindent\textbf{Step 3：}根据【终止条件、筛选准则或评价指标】确定最终结果，并保存【问题二需要调用的输出】。

求解使用【软件、求解器或程序语言】完成，关键参数取值及其来源为【参数设置与依据】，完整程序见附录【编号】。

\subsubsection{求解结果与分析}

求解得到【真实的核心结果】。由【图或表编号】可知，【描述主要现象或变化规律】；
其中【关键结果】说明【数学含义或现实含义】。结合题目要求，问题一的结论为【直接回答问题一】，
并将【真实的中间结果】传递至问题二。

\subsubsection{模型检验与灵敏度分析}

针对【本问模型类型与主要结论】，选取【检验方法】进行验证。
检验指标【指标名称】的结果为【真实结果】，与【判定标准或基准方法】对比后表明【检验结论】。
若【关键参数】具有现实不确定性且会影响主要结论，则在【有依据的变化范围】内进行扰动，
分析【输出数值或决策方案】的变化并说明【稳定区间、临界点或敏感方向】；若无必要，可删除本句而不强行开展灵敏度分析。

%%%%%%%%%%%%%%%%%%%%%%%%%%%%%%%%%%%%%%%%%%%%%%%%%%%%%%%%%%%%%
% ==================== 5.2：问题二的模型建立与求解 ====================
% ✓ 自查：问题一输出已真实进入本问；新增条件已体现；比较标准明确；检验不过度展开。

\subsection{问题二的模型建立与求解}

\subsubsection{建模思路与模型选择}

问题二在问题一的基础上进一步要求【概括问题二的任务】，新增了【新增条件或研究要求】。
本问以问题一得到的【参数、特征或中间结果】作为【本问中的作用】，重点解决【核心数学问题】。
根据【新增条件和数据特点】，选择【主模型名称】对问题一的方法进行【扩展、修正或重新建模】，
并采用【必要的对比模型或基准方案】从【比较指标】方面评价改进效果。
如需额外处理数据，则说明【处理对象、方法及理由】，不单独设置无实际内容的数据预处理小节。

\subsubsection{模型的建立}

在问题一符号体系的基础上，新增【变量或参数】表示【含义】，其余符号与前文保持一致。
依据【理论、机理或题目规则】建立本问的【核心方程、目标函数或判别规则】，并加入【新增约束或条件】：

\begin{equation}
  \text{【问题二核心方程、目标函数或判别规则】},
  \label{eq:q2_core}
\end{equation}

\begin{equation}
  \text{【问题二新增约束、边界条件或适用条件】}.
  \label{eq:q2_constraint}
\end{equation}

\noindent\textbf{模型汇总：}
综合问题一的输入、本问新增条件和变量范围，得到问题二的完整模型

\begin{equation}
  \mathcal{M}_2:\quad
  \left\{
  \begin{aligned}
    &\text{【核心模型表达式】},\\
    &\text{【继承条件与新增约束】},\\
    &\text{【变量范围、初始条件或边界条件】}.
  \end{aligned}
  \right.
  \label{eq:q2_summary}
\end{equation}

\subsubsection{模型的求解}

\noindent\textbf{Step 1：}读取问题一输出的【中间结果】并结合【新增数据或条件】，构造问题二的求解输入。

\noindent\textbf{Step 2：}采用【算法或计算方法】求解【主模型】，记录【目标值、状态量或评价指标】。

\noindent\textbf{Step 3：}在相同输入和评价口径下求解【对比模型或基准方案】，根据【比较指标】确定最终结果并输出【传递给问题三的内容】。

求解使用【软件、求解器或程序语言】完成，关键参数和终止条件为【真实设置及依据】，完整程序见附录【编号】。

\subsubsection{求解结果与分析}

问题二得到【真实的核心结果】。由【图或表编号】可知，【主模型】在【比较指标】方面表现为【真实对比结果】，
说明新增的【条件、机制或改进措施】对【研究对象】产生了【实际影响】。
据此，问题二的答案为【直接结论】，并将【参数、方案或预测结果】用于问题三。

\subsubsection{模型检验与灵敏度分析}

围绕【本问主要结论】，采用【与模型类型匹配的检验方法】验证【拟合、排序、约束或机制】是否可靠。
指标【指标名称】及其真实结果为【结果】，依据【判定标准或基准】可得【检验结论】。
仅当【关键参数】存在现实波动且可能改变结论时，分析其在【有依据的范围】内变化对
【核心输出或最优策略】的影响，并报告【稳定区间、敏感方向或策略切换点】。

%%%%%%%%%%%%%%%%%%%%%%%%%%%%%%%%%%%%%%%%%%%%%%%%%%%%%%%%%%%%%
% ==================== 5.3：问题三的模型建立与求解 ====================
% ✓ 自查：前两问结果与本问条件衔接清楚；模型完整；步骤可复现；结论由真实结果支持。

\subsection{问题三的模型建立与求解}

\subsubsection{建模思路与模型选择}

问题三要求在前两问结果的基础上完成【概括问题三的任务】，需要同时考虑【关键条件一】与【关键条件二】。
本问将【问题一结果】和【问题二结果】分别作为【参数、输入或约束】，其核心是【关键数学问题】。
根据【问题结构、变量性质或求解规模】，选择【主模型名称】描述【目标与约束之间的关系】，
必要时以【基准方法或对比模型】检验所选方案的有效性。
若新增数据必须处理，则在此说明【处理方法及理由】，不固定设置数据预处理标题。

\subsubsection{模型的建立}

定义【本问新增变量】为【符号及含义】，并沿用前两问中的【已有变量或参数】。
依据【现实机理、数学原理或题目规则】建立【核心方程、目标函数或状态转移关系】，同时满足【主要约束】：

\begin{equation}
  \text{【问题三核心方程、目标函数或状态转移关系】},
  \label{eq:q3_core}
\end{equation}

\begin{equation}
  \text{【问题三约束、边界条件或适用条件】}.
  \label{eq:q3_constraint}
\end{equation}

\noindent\textbf{模型汇总：}
汇总前两问输入、本问目标和全部约束，得到问题三的完整模型

\begin{equation}
  \mathcal{M}_3:\quad
  \left\{
  \begin{aligned}
    &\text{【核心模型表达式】},\\
    &\text{【全部主要约束或演化规则】},\\
    &\text{【变量范围、初始条件或边界条件】}.
  \end{aligned}
  \right.
  \label{eq:q3_summary}
\end{equation}

\subsubsection{模型的求解}

\noindent\textbf{Step 1：}整合前两问输出与【本问新增条件】，形成统一的【参数集合、状态初值或候选方案集】。

\noindent\textbf{Step 2：}采用【算法或计算方法】依次完成【关键计算过程】，并记录【必要的中间结果】。

\noindent\textbf{Step 3：}依据【目标、终止条件或筛选准则】获得最终【预测、方案或机理参数】，并输出问题四所需的【结果】。

求解使用【软件、求解器或程序语言】完成，关键参数、初值和终止条件为【真实设置及依据】，完整程序见附录【编号】。

\subsubsection{求解结果与分析}

求解结果为【真实结果】。由【图或表编号】可见，【描述核心规律、方案特征或变化趋势】；
进一步比较【对象或情景】可知【差异及原因】。因此，问题三的结论为【直接回答问题三】，
其中【必要输出】将作为问题四综合分析的依据。

\subsubsection{模型检验与灵敏度分析}

根据【本问模型类型】，采用【检验方法】检查【关键数学关系、约束可行性或结果精度】。
真实检验结果【指标或现象】达到【判定标准或基准】，说明【检验结论】。
若【核心不确定参数】可能影响最终方案，则在【现实变化范围】内进行扰动，重点分析
【数值变化是否引起决策结构变化】，并报告【稳定区间、临界阈值或敏感方向】。

%%%%%%%%%%%%%%%%%%%%%%%%%%%%%%%%%%%%%%%%%%%%%%%%%%%%%%%%%%%%%
% ==================== 5.4：问题四的模型建立与求解 ====================
% ✓ 自查：综合使用前三问成果；最终结论逐项回应题目；检验服务于结论而非形式完整。

\subsection{问题四的模型建立与求解}

\subsubsection{建模思路与模型选择}

问题四要求综合前三问结果完成【最终预测、决策、优化或机理解释任务】。
本问以【问题一结果】、【问题二结果】和【问题三结果】作为【输入、参数或备选方案】，
同时考虑【最终新增条件或现实限制】，需要回答【核心问题】。
据此选择【主模型名称或综合分析方法】完成【最终任务】，并以【评价标准或基准方案】衡量结果质量。
如有必要的数据整理，只说明与最终模型直接相关的【处理内容与理由】。

\subsubsection{模型的建立}

定义【最终决策变量、输出变量或状态变量】为【符号及含义】，定义【新增参数】为【符号及含义】。
结合前三问结论和【现实约束】，建立【核心方程、目标函数、决策规则或机理关系】：

\begin{equation}
  \text{【问题四核心方程、目标函数、决策规则或机理关系】},
  \label{eq:q4_core}
\end{equation}

\begin{equation}
  \text{【问题四约束、边界条件或适用条件】}.
  \label{eq:q4_constraint}
\end{equation}

\noindent\textbf{模型汇总：}
综合全部输入、目标、约束和变量范围，得到问题四的完整模型

\begin{equation}
  \mathcal{M}_4:\quad
  \left\{
  \begin{aligned}
    &\text{【最终模型表达式】},\\
    &\text{【全部主要约束或决策规则】},\\
    &\text{【变量范围、初始条件或边界条件】}.
  \end{aligned}
  \right.
  \label{eq:q4_summary}
\end{equation}

\subsubsection{模型的求解}

\noindent\textbf{Step 1：}汇总前三问的【真实输出】和问题四的【新增输入】，统一符号、单位与评价口径。

\noindent\textbf{Step 2：}采用【算法、求解器或综合分析流程】计算【最终预测、决策变量或候选方案】。

\noindent\textbf{Step 3：}根据【评价标准、可行性条件或决策准则】筛选最终结果，并形成【题目要求的输出形式】。

求解使用【软件、求解器或程序语言】完成，关键参数与终止条件为【真实设置及依据】，完整程序见附录【编号】。

\subsubsection{求解结果与分析}

最终得到【真实的预测值、最优方案、决策建议或机理结论】。
由【图或表编号】可知，【主要规律或方案特征】，其产生原因是【结合模型与现实条件解释原因】。
上述结果分别回答了【题目要求一】与【题目要求二】，最终结论为【完整回答问题四】。

\subsubsection{模型检验与灵敏度分析}

围绕最终结论，选择【与主模型匹配的检验方法】验证【结果精度、约束可行性、稳定性或机理合理性】。
检验得到【真实指标、对照结果或边界表现】，依据【判断标准】可得【检验结论】。
若存在会改变最终决策的【现实不确定参数】，则在【有依据的范围】内分析其影响，
说明【主要趋势、稳定区间或策略切换点】；不得用程序运行正常代替数学意义上的模型检验。

